\section{The functional analysis forced on us by the calculation}
\label{sec:functional-operator-algebra}
% BEGIN CURATED SUBJECT INDEX
\index{unbounded operator}
\index{topology}
\index{measure theory}
\index{resolvent}
\index{direct integral}
\index{von Neumann algebra@\vN\ algebra}
% END CURATED SUBJECT INDEX

The elementary formulas of the preceding section already demand most of the
functional analysis used in scattering theory.  The strong limit in the wave
operator, the boundary value of the resolvent, the delta-normalized
eigenfunction, and the analytic continuation of a matrix element are four
different limiting operations.  Treating them as if they were ordinary limits
of finite matrices hides both the power of the theory and the origin of many
apparent paradoxes.  This section develops the relevant language from first
principles and keeps a physical example next to every abstract definition.
The order is deliberately diagnostic rather than encyclopedic: a topology is
introduced to specify a limit, an operator domain to specify an observable, a
direct integral to specify the continuum, and a \vN\ algebra to specify
which decompositions and sectors are measurable.
Standard references include Refs.~\cite{ReedSimonI:1980,Takesaki:2002,
vonNeumann:1949Reduction,Yafaev:1992,Haag:1996}.

\subsection{Topology decides what convergence means}
\label{subsec:topology-convergence}
% BEGIN CURATED SUBJECT INDEX
\index{operator!adjoint and self-adjoint}
\index{Hilbert space}
\index{locally convex space}
\index{topology}
\index{measure theory}
\index{resolvent}
\index{von Neumann algebra@\vN\ algebra}
\index{trace-class operator}
\index{wave operator}
% END CURATED SUBJECT INDEX

A topology on a set $X$ is a collection of subsets declared to be open, closed
under arbitrary unions and finite intersections and containing both $X$ and
the empty set.  It determines continuity and convergence before a distance is
chosen.  A metric gives a topology, but many spaces of operators and
distributions carry useful topologies that are not naturally described by a
single metric.  On a vector space, addition and scalar multiplication are
required to be continuous.  A locally convex topology is generated by a
family of seminorms $p_\alpha$, where a seminorm obeys the triangle inequality
and homogeneity but may vanish on nonzero vectors.

The norm topology of a Hilbert space $\Hil$ is generated by
$p(\psi)=\|\psi\|$.  A sequence converges weakly if all transition
amplitudes against fixed test vectors converge:
\begin{equation}
  \psi_n\rightharpoonup\psi
  \quad\Longleftrightarrow\quad
  \dualpair{\phi}{\psi_n}
  \to\dualpair{\phi}{\psi}
  \quad\text{for every $\phi\in\Hil$} .
  \label{eq:weak-vector-convergence}
\end{equation}
Norm convergence implies weak convergence, but the converse fails.  An
orthonormal sequence is the standard example: $e_n\rightharpoonup0$ although
$\|e_n\|=1$.  Physically, probability can escape to infinity while every
fixed localized detector sees a vanishing amplitude.  This is exactly the
kind of limit that occurs in scattering.

For bounded operators $A_n\in\BHil$, the most frequently used topologies are
summarized in Table~\ref{tab:operator-topologies}.  The adjective
\textit{strong} refers to norm convergence after acting on each vector; it
does not mean operator-norm convergence.  On an infinite-dimensional space
the distinctions are strict.  Nets rather than sequences are required for
complete generality, because weak operator topologies need not be first
countable.  For a separable Hilbert space, however, the strong and weak
operator topologies are metrizable on norm-bounded sets, which covers many
physics applications.

\begin{table}[t]
  \centering
  \small
  \begin{tabular}{>{\raggedright\arraybackslash}p{0.19\textwidth}
                  >{\raggedright\arraybackslash}p{0.34\textwidth}
                  >{\raggedright\arraybackslash}p{0.35\textwidth}}
    \toprule
    topology & convergence of $A_n$ to $A$ & typical role in physics \\
    \midrule
    operator norm &
    $\|A_n-A\|\to0$ &
    uniform control on every unit vector; stable functional calculus \\
\addlinespace
    strong operator &
    $\|(A_n-A)\psi\|\to0$ for every $\psi$ &
    wave operators, unitary dynamics, and strong sums of projections \\
\addlinespace
    weak operator &
    $\dualpair{\phi}{(A_n-A)\psi}\to0$ for every
    $\phi$, $\psi$ &
    convergence of all fixed transition amplitudes \\
\addlinespace
    ultraweak &
    convergence against every trace-class functional &
    normal states and \vN\ algebras \\
\addlinespace
    strong resolvent &
    $(A_n-z)^{-1}\psi\to(A-z)^{-1}\psi$ for one, hence all,
    $z\in\complex\setminus\real$ &
    limits of unbounded self-adjoint Hamiltonians \\
    \bottomrule
  \end{tabular}
  \caption{Common operator topologies.  The first four rows concern bounded
  operators.  Strong-resolvent convergence is a topology adapted to
  unbounded self-adjoint operators.}
  \label{tab:operator-topologies}
\end{table}

Topology and measure are related but distinct.  A topology on $X$ generates a
Borel sigma-algebra $\Sigma$, the smallest sigma-algebra containing the open
sets.  A measure space $(X,\Sigma,\mu)$ then specifies which subsets are
measurable and how they are weighted.  Spectral theory uses topology to discuss
continuity in energy and measure theory to identify functions that differ only
on sets of $\mu$ measure zero.  The phrase ``at almost every energy'' is
therefore not a technical evasion; it is part of the identification defining
an $L^2$ state.

\subsection{Unbounded operators: the domain is part of the observable}
\label{subsec:unbounded-operators}
% BEGIN CURATED SUBJECT INDEX
\index{operator!adjoint and self-adjoint}
\index{self-adjoint extension}
\index{unbounded operator}
\index{operator!closed}
\index{resolvent}
% END CURATED SUBJECT INDEX

A differential Hamiltonian is not merely a formula $H\psi$.  It is a pair
consisting of that action and a domain $D(H)\subset\Hil$.  The graph
\begin{equation}
  \mathcal G(H)=\{(\psi,H\psi):\psi\in D(H)\}
  \subset\Hil\oplus\Hil
  \label{eq:operator-graph}
\end{equation}
is closed if convergence of both $\psi_n$ and $H\psi_n$ implies that the
limit lies in the graph.  An operator is closable if the closure of its graph
is the graph of an operator.  Closedness is the correct replacement for
boundedness in much of spectral theory.

For a densely defined operator $A$, the adjoint domain is
\begin{equation}
  D(A^\dagger)
  =\{\phi\in\Hil:
  \text{$\psi\mapsto\dualpair{\phi}{A\psi}$ is norm continuous on $D(A)$}\} .
  \label{eq:adjoint-domain}
\end{equation}
If $A\subset A^\dagger$, it is symmetric; if $A=A^\dagger$, including
equality of domains, it is self-adjoint.  A symmetric radial differential
expression can acquire different self-adjoint extensions from different
boundary conditions at a singular endpoint.  The deficiency spaces
$\ker(A^\dagger\mp\rmi)$ count the missing boundary data.  This is why
``Hermitian differential expression'' is not a sufficient specification of a
quantum Hamiltonian.

For a self-adjoint $H$, the resolvent $(z-H)^{-1}$ is bounded for
$z\notin\real$ and satisfies
\begin{equation}
  \|(z-H)^{-1}\|\leq\frac{1}{|\im z|} .
  \label{eq:self-adjoint-resolvent-bound}
\end{equation}
Thus an unbounded Hamiltonian can be studied through bounded resolvents,
unitaries $\rme^{-\rmi tH}$, and spectral projections.  Operator algebras
usually contain these bounded functions of $H$, while $H$ itself is described
as an operator \textit{affiliated} with the algebra.  Quadratic-form methods
provide another route: a closed semibounded form determines a unique
self-adjoint operator and permits singular potentials for which the naive
operator sum is inconvenient \cite{ReedSimonI:1980}.

\subsection{Projection-valued measures and the spectral theorem}
\label{subsec:pvm-spectral-theorem}
% BEGIN CURATED SUBJECT INDEX
\index{operator!adjoint and self-adjoint}
\index{topology}
\index{spectral theorem}
\index{spectral measure}
\index{spectrum!point}
\index{resolvent}
\index{GNS construction}
\index{wave packet}
% END CURATED SUBJECT INDEX

A projection-valued measure $P$ assigns an orthogonal projection $P(\Delta)$
to each Borel set $\Delta\subset\real$, with
$P(\real)=\identity$ and countable additivity in the strong operator
topology.  The spectral theorem states that a self-adjoint operator can be
written
\begin{equation}
  H=\int_{\real}E\dd P(E),
  \qquad
  f(H)=\int_{\real}f(E)\dd P(E) .
  \label{eq:pvm-spectral-theorem}
\end{equation}
The domain of $H$ is not suppressed by this notation:
\begin{equation}
  D(H)=\left\{\psi\in\Hil:
  \int_{\real}E^2\dd\mu_\psi(E)<\infty\right\},
  \qquad
  \mu_\psi(\Delta)
  =\dualpair{\psi}{P(\Delta)\psi} .
  \label{eq:spectral-measure-vector}
\end{equation}
The scalar measure $\mu_\psi$ is the energy distribution of the state.
Its pure-point, absolutely continuous, and singular-continuous parts induce
the corresponding orthogonal decomposition of $\Hil$.  A bound state gives
an atom of the measure.  A scattering wave packet has an absolutely
continuous density.  A delta-normalized ``eigenket'' at one continuum energy
is not a vector in $\Hil$; it is a distributional coordinate for this
measure decomposition.

The spectral theorem also places a sharp restriction on resonances.  Every
spectral projection of a self-adjoint $H$ lies on the real spectrum.  A pole at
$\resonantE\notin\real$ cannot be added to
Eq.~\eqref{eq:pvm-spectral-theorem} as another orthogonal projector.  It
appears only after continuing suitably tested resolvent matrix elements, or as
a genuine eigenvalue of a different, non-self-adjoint realization such as
$H_\theta$.  This boundary between spectral measure and analytic
continuation is the mathematical core of resonance theory.

\subsection{Direct integrals: a continuum of channel Hilbert spaces}
\label{subsec:direct-integrals}
% BEGIN CURATED SUBJECT INDEX
\index{Banach space}
\index{Hilbert space}
\index{measure theory}
\index{spectral theorem}
\index{direct integral}
\index{decomposable operator}
\index{commutant}
% END CURATED SUBJECT INDEX

A direct sum describes discrete alternatives.  A direct integral is its
measure-theoretic analogue.  Let $(X,\Sigma,\mu)$ be a sigma-finite standard
measure space.  A field $x\mapsto\Hil_x$ of separable Hilbert spaces is called
\term{measurable} after choosing a countable fundamental family of sections
$e_n(x)\in\Hil_x$ such that
\begin{equation}
  \begin{aligned}
    x&\longmapsto\dualpair{e_m(x)}{e_n(x)}_{\Hil_x}
      &&\text{is measurable},\\
    \overline{\operatorname{span}}\{e_n(x)\}&=\Hil_x
      &&\text{for almost every $x$} .
  \end{aligned}
  \label{eq:measurable-field-fundamental-family}
\end{equation}
A section $\psi(x)$ is measurable when all scalar functions
$x\mapsto\dualpair{e_n(x)}{\psi(x)}$ are measurable.  Different
fundamental families defining the same measurable sections give the same
measurable field.  This condition is essential: there is no useful direct
integral of an arbitrary uncoordinated collection of Hilbert spaces.

The direct integral
\begin{equation}
  \Hil=\int_X^\oplus\Hil_x\dd\mu(x)
  \label{eq:direct-integral-Hilbert}
\end{equation}
consists of measurable sections $x\mapsto\psi(x)$, identified when they
agree almost everywhere, for which
\begin{equation}
  \|\psi\|^2
  =\int_X\|\psi(x)\|_{\Hil_x}^2\dd\mu(x)<\infty .
  \label{eq:direct-integral-norm}
\end{equation}
The fiber dimension may vary with $x$.  In scattering, this is exactly what
happens when a new channel opens at a threshold.

A measurable essentially bounded family
$x\mapsto A(x)\in\Banach(\Hil_x)$ defines a decomposable operator
\begin{equation}
  A=\int_X^\oplus A(x)\dd\mu(x),
  \qquad
  (A\psi)(x)=A(x)\psi(x),
  \label{eq:decomposable-operator}
\end{equation}
with $\|A\|=\esssup_x\|A(x)\|$.  The scalar
diagonal algebra
\begin{equation}
  \mathcal D
  =\left\{\int_X^\oplus f(x)\identity_x\dd\mu(x):
  f\in L^\infty(X,\mu)\right\}
  \label{eq:diagonal-algebra}
\end{equation}
is abelian.  For any set $\mathcal X\subset\BHil$, its
\term{commutant} is
\begin{equation}
  \mathcal X'
  =\{B\in\BHil:
  \text{$BA=AB$ for every $A\in\mathcal X$}\}.
  \label{eq:commutant-first-definition}
\end{equation}
Under the standard hypotheses, the commutant $\mathcal D'$ is
the algebra of all decomposable bounded operators.  Thus ``does not mix the
base variable $x$'' is equivalent to ``commutes with every bounded measurable
function of $x$.''

\begin{derivationbox}{Why the multiplicity-one commutant is diagonal}
First suppose $\mu(X)<\infty$ and the fiber is $\complex$.  If $A$ commutes
with every multiplication operator $M_f$, it commutes in particular with
$M_{\chi_\Delta}$ for every measurable $\Delta$.  Therefore $A$ preserves
functions supported in $\Delta$.  Put $g=A\identity$.  For a simple function
$s=\sum_jc_j\chi_{\Delta_j}$,
\begin{equation}
  As=\sum_jc_jA M_{\chi_{\Delta_j}}\identity
  =\sum_jc_jM_{\chi_{\Delta_j}}A\identity=gs.
  \label{eq:multiplicity-one-simple-proof}
\end{equation}
Density of simple functions gives $A=M_g$.  Testing on characteristic
functions of sets where $|g|>\|A\|+\epsilon$ shows that
$g\in L^\infty$ and $\|g\|_\infty\leq\|A\|$.  On a sigma-finite space, apply
the argument on an increasing finite-measure partition; the local multipliers
agree on overlaps.  For an $m(x)$-dimensional fiber, the same support argument
leaves an essentially bounded measurable matrix $A(x)$ acting inside each
fiber.  This is the decomposable-operator statement used below.
\end{derivationbox}

The multiplicity form of the spectral theorem supplies a measure space and a
unitary map such that
\begin{align}
  U\Hil
  &=\int_{\spectrum(H)}^\oplus\mathcal K_E\dd\mu(E),
  \label{eq:spectral-direct-integral}\\
  (UHU^\dagger\psi)(E)
  &=E\psi(E),
  \label{eq:H-multiplication-direct-integral}
\end{align}
where $\dim\mathcal K_E$ is the spectral multiplicity.  The measure and
fibers are unique only up to null sets and measure-class equivalence.  This is
why changing continuum normalization changes density factors but not the
unitary content of the decomposition.

\subsection{Worked example: the free-particle energy representation}
\label{subsec:free-direct-integral}
% BEGIN CURATED SUBJECT INDEX
\index{wave operator}
\index{partial-wave expansion}
\index{scattering!on-shell}
% END CURATED SUBJECT INDEX

The abstract construction becomes elementary in momentum space.  For a
spinless free particle in three dimensions,
\begin{equation}
  \Hil=L^2(\real^3,\rmd^3p),
  \qquad
  (H_0\hat\psi)(\bm p)=\frac{p^2}{2\mu}\hat\psi(\bm p) .
  \label{eq:free-H-momentum}
\end{equation}
Set $p(E)=\sqrt{2\mu E}$.  Since
\begin{equation}
  \rmd^3p=p^2\dd p\dd\Omega
  =\mu p(E)\dd E\dd\Omega,
  \label{eq:free-energy-measure}
\end{equation}
the map
\begin{equation}
  (U_0\hat\psi)(E,\hat{\bm p})
  =[\mu p(E)]^{1/2}
  \hat\psi(p(E)\hat{\bm p})
  \label{eq:free-energy-unitary}
\end{equation}
is unitary from $L^2(\real^3,\rmd^3p)$ onto
\begin{equation}
  \int_0^\infty{}^\oplus
  L^2(S^2,\rmd\Omega)\dd E .
  \label{eq:free-particle-direct-integral}
\end{equation}
The factor $[\mu p(E)]^{1/2}$ is the square root of the density of states
needed to preserve the norm.  In this representation $H_0$ is multiplication
by $E$, while each fiber contains the angular directions at that energy.
Spherical harmonics further decompose the fiber into partial waves.  Spin,
internal states, and reaction channels simply enlarge the fiber.

At this point we have only diagonalized $H_0$; no scattering operator has
been used.  After Stage II constructs the wave operators, the same
representation will turn strong commutation with $H_0$ into an
almost-everywhere family of on-shell operators.  That implication and its
hypotheses are derived in Section~\ref{sec:scattering-direct-integral}.

\subsection{Operator algebras, \vN\ algebras, and states}
\label{subsec:cstar-von-neumann}
% BEGIN CURATED SUBJECT INDEX
\index{operator!adjoint and self-adjoint}
\index{Hilbert space}
\index{topology}
\index{commutant}
\index{von Neumann algebra@\vN\ algebra}
\index{GNS construction}
\index{infrared problem}
% END CURATED SUBJECT INDEX

An operator algebra separates the abstract observables from a particular
choice of vectors.  A $C^*$ algebra $\Alg$ is a norm-complete involutive
algebra satisfying
\begin{equation}
  \|A^\dagger A\|=\|A\|^2 .
  \label{eq:cstar-identity}
\end{equation}
When represented on $\Hil$, it is a norm-closed $*$ subalgebra of
$\BHil$.  A \vN\ algebra $\vNAlg\subset\BHil$ is a unital
$*$ algebra closed in the weak operator topology.  The bicommutant theorem
gives the equivalent characterization
\begin{equation}
  \vNAlg=\vNAlg'' .
  \label{eq:bicommutant-theorem}
\end{equation}
Norm closure is appropriate for uniform approximation of observables; weak
closure also includes limits determined by all matrix elements and produces
the spectral projections of self-adjoint generators.  For this reason von
Neumann algebras are particularly natural for spectral theory and statistical
states \cite{Takesaki:2002,BratteliRobinsonI:1987}.

An algebraic state is a positive normalized linear functional
\begin{equation}
  \omega(A^\dagger A)\geq0,
  \qquad
  \omega(\identity)=1 .
  \label{eq:algebraic-state}
\end{equation}
The Gelfand--Naimark--Segal (GNS) construction associates to it a Hilbert space
$\Hil_\omega$, a representation $\pi_\omega(\Alg)$, and a cyclic
vector $\Omega_\omega$ such that
\begin{equation}
  \omega(A)
  =\dualpair{\Omega_\omega}
  {\pi_\omega(A)\Omega_\omega} .
  \label{eq:GNS-state}
\end{equation}
This reverses the usual order of presentation: the state and observable
algebra determine a representation, rather than every physical state being
assumed in advance to lie in one preferred Fock or Schr\"odinger Hilbert space.
For finitely many degrees of freedom the familiar irreducible representations
are essentially unique under regularity hypotheses.  For infinitely many
degrees of freedom, inequivalent representations are generic.  The infrared
problem of QED in Section~\ref{sec:long-range-ir} is a concrete scattering
reason that this distinction becomes unavoidable.

A state on a \vN\ algebra is normal when it is continuous under
increasing strong limits of projections, equivalently when it belongs to the
predual $\vNAlg_*$.  On $\BHil$, normal states have density matrices,
$\omega(A)=\tr(\rho A)$.  Singular states need not.  Unbounded observables
are not elements of $\vNAlg$, but their spectral projections may lie in it;
one then says that the operator is affiliated with $\vNAlg$.

\subsection{Abelian spectral algebras and central decomposition}
\label{subsec:central-decomposition}
% BEGIN CURATED SUBJECT INDEX
\index{operator!adjoint and self-adjoint}
\index{Banach space}
\index{Hilbert space}
\index{measure theory}
\index{direct integral}
\index{direct integral!fiber}
\index{decomposable operator}
\index{commutant}
\index{von Neumann algebra@\vN\ algebra}
\index{central decomposition}
% END CURATED SUBJECT INDEX

The bounded Borel functions of a self-adjoint $H$ generate an abelian von
Neumann algebra
\begin{equation}
  W^*(H)=\{f(H):\text{$f$ bounded and Borel}\}'' .
  \label{eq:von-neumann-generated-H}
\end{equation}
In the spectral representation it is the diagonal algebra
$L^\infty(X,\mu)$ acting as $f(E)\identity_E$.  Its commutant consists
of the decomposable operators acting inside the multiplicity fibers.  Stage II
will show that strong commutation with $H_0$ is therefore equivalent to
membership in $W^*(H_0)'$.

There is a broader direct-integral theorem.  Let $\vNAlg$ act on a separable
Hilbert space and assume the usual countable-decomposability hypotheses.  Its
center
\begin{equation}
  \Center(\vNAlg)=\vNAlg\cap\vNAlg'
  \label{eq:center-von-neumann}
\end{equation}
is abelian and can be represented as $L^\infty(X,\mu)$ on a standard
measure space.  The reduction theory of \vN\ then gives
\begin{align}
  \Hil&=\int_X^\oplus\Hil_x\dd\mu(x),
  \label{eq:central-H-decomposition}\\
  \vNAlg&=\int_X^\oplus\vNAlg_x\dd\mu(x),
  \qquad
  \Center(\vNAlg_x)=\complex\identity_x
  \quad\text{for almost every $x$} .
  \label{eq:central-algebra-decomposition}
\end{align}
The fibers $\vNAlg_x$ are factors \cite{vonNeumann:1949Reduction,
Takesaki:2002}.  In finite dimensions this is the familiar decomposition of a
block-diagonal algebra into full matrix blocks.  The direct integral is the
continuous version.  This theorem decomposes a \textit{specified
representation} of one \vN\ algebra.  It does not by itself decide
which representations are physically admissible, nor does it classify all
superselection sectors of an abstract observable algebra.

Factors are classified coarsely into types.  Type I factors contain minimal
projections and reproduce ordinary quantum mechanics, with
$\vNAlg_x\simeq\Banach(\Hil_x)$.  Type II factors have no minimal
projections but admit a faithful semifinite trace.  Type III factors admit no
nonzero finite projections and no ordinary trace; local observable algebras in
relativistic QFT are typically type III under standard assumptions
\cite{Haag:1996}.  This explains why a local QFT region does not behave like
a finite tensor factor with an intrinsic density matrix, even though detector
probabilities remain well defined.

Three notions should not be conflated.  The spectral direct integral of
$H$ is taken over energy and diagonalizes $W^*(H)$.  The central decomposition
of a represented observable algebra is taken over the center and resolves
that representation into factors.  A superselection analysis additionally
selects and compares physically admissible representations of the abstract
algebra, often by a localization or asymptotic criterion.  The base of a
central decomposition represents sector labels only when the chosen
representation was constructed to contain such a mixture and its center
records those labels.  The three constructions use related measurable
technology, but energy need not be a superselection variable: generic
observables can change it.  In infrared QED, by contrast, asymptotic electric
flux or soft charge can label disjoint representations.  A statistical
mixture of them may admit a direct integral, but that operation creates no
unitary intertwiner between disjoint fibers.

\subsection{Locally convex test spaces and generalized eigenvectors}
\label{subsec:locally-convex-rhs-preview}
% BEGIN CURATED SUBJECT INDEX
\index{operator!adjoint and self-adjoint}
\index{locally convex space}
\index{topology}
\index{spectral theorem}
\index{complex scaling}
\index{rigged Hilbert space}
% END CURATED SUBJECT INDEX

The Hilbert norm deliberately forgets pointwise values, derivatives, and
analytic continuation.  Scattering theory often needs all three.  Choose a
dense invariant space $\PhiSpace\subset\Hil$ and give it a locally convex
topology generated, for example, by seminorms
\begin{equation}
  p_{m,n}(\phi)
  =\left\|\braket{x}^m
  (1+H)^n\phi\right\|,
  \qquad \text{$m$, $n\in\naturalnumbers$} .
  \label{eq:test-space-seminorms}
\end{equation}
For the free particle, the Schwartz space is the canonical model.  Its
topology controls every polynomially weighted derivative, and its continuous
dual contains tempered distributions such as plane waves and Dirac deltas.
Analytic or Hardy seminorms can be added when contour deformation is required.

Continuity of a functional $F$ on $\PhiSpace$ means that for some finite set
of seminorms and a constant $C$,
\begin{equation}
  |F(\phi)|\leq C\sum_{j=1}^N p_j(\phi) .
  \label{eq:functional-continuity-seminorm}
\end{equation}
This gives the Gelfand triple
$\PhiSpace\subset\Hil\subset\PhiSpace^\times$ used in
Section~\ref{sec:rigged-hilbert}.  A generalized eigenvector
$\ket{E,\alpha}\in\PhiSpace^\times$ satisfies
\begin{equation}
  \dualpair{H\phi}{E,\alpha}
  =E\dualpair{\phi}{E,\alpha},
  \qquad \phi\in\PhiSpace,
  \label{eq:generalized-eigenvector-preview}
\end{equation}
where $\alpha$ labels the spectral fiber.  Under nuclearity and invariance
hypotheses, the nuclear spectral theorem turns the direct-integral expansion
into a generalized-eigenvector expansion \cite{Gelfand:1964}.  The topology
is therefore the bridge between an almost-everywhere spectral representation
and the physicist's energy kets.

The choice of $\PhiSpace$ is not arbitrary bookkeeping.  A topology too weak
makes evaluation or continuation discontinuous; one too strong can exclude
physically prepared packets.  Hardy spaces encode a time-oriented analytic
boundary value, Schwartz spaces encode rapid localization, and dilation-
analytic vectors encode complex scaling.  Different choices can reveal
different generalized spectra without changing the self-adjoint spectrum in
$\Hil$.

\subsection{Riesz projections and what an operator algebra does not add}
\label{subsec:riesz-operator-algebra}
% BEGIN CURATED SUBJECT INDEX
\index{operator!adjoint and self-adjoint}
\index{topology}
\index{measure theory}
\index{spectral measure}
\index{resolvent}
\index{resolvent set}
\index{Riesz projection}
\index{direct integral}
\index{direct integral!fiber}
\index{decomposable operator}
% END CURATED SUBJECT INDEX

For a closed, not necessarily normal operator $K$, an isolated spectral set
inside a contour $\Gamma$ has the Riesz projection
\begin{equation}
  P_\Gamma(K)
  =\frac{1}{2\pi\rmi}\oint_\Gamma
  (z-K)^{-1}\dd z .
  \label{eq:riesz-projection-general}
\end{equation}
It is idempotent but need not be orthogonal.  Its range is the generalized
eigenspace, including Jordan vectors.  For a complex-scaled Hamiltonian, an
exposed resonance has precisely such a finite-rank projection.  At an
exceptional point the rank of the projection can remain fixed while the
geometric eigenspace shrinks and a nilpotent Jordan part develops.

For the original self-adjoint $H$, however, a resonance behind the cut is not
an isolated part of $\spectrum(H)$, so Eq.~\eqref{eq:riesz-projection-general}
cannot encircle it on the physical resolvent set.  Neither norm closure nor
weak closure of the observable algebra manufactures the missing analytic
sheet.  One must supply analytic vectors, weighted spaces, a Jost construction,
or a complex deformation.  Operator algebras organize observables,
representations, and measurable decompositions; resonance theory adds the
analytic continuation.  Their interface, rather than an identification of the
two, is the useful unifying structure.

\begin{table}[t]
  \centering
  \small
  \begin{tabular}{>{\raggedright\arraybackslash}p{0.30\textwidth}
                  >{\raggedright\arraybackslash}p{0.55\textwidth}}
    \toprule
    physical phrase & precise mathematical object \\
    \midrule
    energy distribution of a packet & scalar spectral measure
    $\mu_\psi(\Delta)=\dualpair{\psi}{P(\Delta)\psi}$ \\
    \addlinespace
    continuum channels at fixed energy & multiplicity fiber
    $\mathcal K_E$ in a direct integral \\
    \addlinespace
    energy-preserving bounded map & decomposable fiber $A(E)$ of an operator
    $A\in W^*(H_0)'$ \\
    \addlinespace
    generalized plane or Gamow wave & continuous functional on a chosen test
    topology \\
    \addlinespace
    isolated complex-scaled resonance & finite-rank Riesz projection of
    $H_\theta$ \\
    \addlinespace
    factor component of a representation & fiber in its central
    decomposition; a separate selection criterion is required before calling
    it a superselection sector \\
    \bottomrule
  \end{tabular}
  \caption{A functional-analytic dictionary for scattering theory.}
  \label{tab:functional-dictionary}
\end{table}
